\section{Conditional face transport and likelihood drift}
\label{sec:tpc44-face-transport}

The coefficient descent tracks the centered spectrum.  There is a
second exact descent that preserves the nonnegativity of \(Z_0\).  For
\(0\le r\le R\), define
\begin{equation}
 A_r:=
 \E_{\varepsilon_{p_1},\ldots,\varepsilon_{p_r}}
 Z_0(\varepsilon_{\le p_r},-\one_{>p_r}).
 \label{eq:tpc44-Ar}
\end{equation}
Then
\begin{equation}
 A_0=Z_0(-\one),
 \qquad A_R=\cD_0,
 \label{eq:tpc44-Ar-endpoints}
\end{equation}
and the reverse Walsh formula gives
\begin{equation}
 A_r=\cD_0+
 \sum_{\substack{\kappa>1\\\Pminus(\kappa)>p_r}}
 \mu(\kappa)C(\kappa).
 \label{eq:tpc44-Ar-spectrum}
\end{equation}
For an arbitrary cutoff \(z\), use \eqref{eq:tpc44-cutoff-index} and write
\begin{equation}
 A_z:=A_{r(z)},
 \qquad E_z:=E_{r(z)}.
 \label{eq:tpc44-Az-definition}
\end{equation}

\begin{proposition}[Exact least--prime increment]
\label{prop:tpc44-least-prime-increment}
Put
\begin{equation}
 L_r:=A_{r-1}-A_r.
 \label{eq:tpc44-Lr}
\end{equation}
Then
\begin{equation}
 L_r=\sum_{\Pminus(\kappa)=p_r}\mu(\kappa)C(\kappa).
 \label{eq:tpc44-Lr-spectrum}
\end{equation}
Moreover,
\begin{equation}
 0\le A_{r-1}\le2A_r.
 \label{eq:tpc44-face-doubling}
\end{equation}
If \(Z_0(-\one)>0\), all \(A_r>0\), and
\begin{equation}
 \delta_r:=\frac{L_r}{A_r}\in[-1,1]
 \label{eq:tpc44-delta}
\end{equation}
satisfies
\begin{equation}
 \boxed{
 \frac{Z_0(-\one)}{\cD_0}
 =\prod_{r=1}^{R}(1+\delta_r).}
 \label{eq:tpc44-face-product}
\end{equation}
\end{proposition}

\begin{proof}
Subtract \eqref{eq:tpc44-Ar-spectrum} at adjacent indices to obtain the
increment formula.  Before averaging \(\varepsilon_{p_r}\), let \(B_-\)
and \(B_+\) be the two nonnegative conditional averages obtained by
setting that sign to \(-1\) and \(+1\), respectively.  Then
\begin{equation}
 A_{r-1}=B_-,
 \qquad A_r=\frac{B_-+B_+}{2},
 \label{eq:tpc44-two-face-average}
\end{equation}
which proves \eqref{eq:tpc44-face-doubling} and
\(-1\le\delta_r\le1\).  If some \(A_r=0\), then \(A_{r-1}=0\), so
positive \(A_0\) forces every \(A_r>0\).  Finally
\(
 A_{r-1}=A_r(1+\delta_r)
\); telescope from \(r=R\) to \(r=1\).
\end{proof}

The increment also has an exact Gram interpretation.  After averaging
the coordinates \(p_1,\ldots,p_{r-1}\), group row labels by their
incidence pattern \(B\) on those coordinates.  With the later signs fixed
at \(-1\), let \(V_{B,0}\) and \(V_{B,1}\) be the resulting coefficient
vectors in the direct sum of the output spaces, according as \(p_r\)
does not or does divide the label.  Orthogonality in the averaged
coordinates gives
\begin{align}
 A_{r-1}
 &=\sum_B\|V_{B,0}-V_{B,1}\|^2,
 \label{eq:tpc44-Gram-face-minus}\\
 A_r
 &=\sum_B\bigl(\|V_{B,0}\|^2+\|V_{B,1}\|^2\bigr),
 \label{eq:tpc44-Gram-face-average}\\
 L_r
 &=-2\operatorname{Re}\sum_B
   \langle V_{B,0},V_{B,1}\rangle.
 \label{eq:tpc44-Gram-face-increment}
\end{align}
Consequently, whenever \(A_r>0\), the drift \(\delta_r=L_r/A_r\) is a
normalized real cross inner product between the \(p_r\)-free and
\(p_r\)-divisible blocks.

The product has a probability interpretation.  Let \(U\) be uniform
measure on the active prime cube and define the energy--biased measure
\begin{equation}
 d\QZ:=\frac{Z_0}{\cD_0}\,dU.
 \label{eq:tpc44-energy-biased-measure}
\end{equation}
For \(0\le r\le R\), let
\begin{equation}
 \mathcal E_r:=
 \{\varepsilon_{p_{r+1}}=\cdots=\varepsilon_{p_R}=-1\}.
 \label{eq:tpc44-tail-face-event}
\end{equation}

\begin{theorem}[Boundary likelihood chain]
\label{thm:tpc44-likelihood-chain}
If \(Z_0(-\one)>0\), then
\begin{equation}
 1+\delta_r
 =2\QZ\!\left(
 \varepsilon_{p_r}=-1\mid\mathcal E_r
 \right),
 \label{eq:tpc44-conditional-likelihood}
\end{equation}
and therefore
\begin{equation}
 \frac{Z_0(-\one)}{\cD_0}
 =\prod_{r=1}^{R}
 2\QZ\!\left(
 \varepsilon_{p_r}=-1
 \mid
 \varepsilon_{p_{r+1}}=\cdots=\varepsilon_{p_R}=-1
 \right).
 \label{eq:tpc44-likelihood-product}
\end{equation}
\end{theorem}

\begin{proof}
The uniform probability of \(\mathcal E_r\) is \(2^{-(R-r)}\), and
conditioning the energy on that face gives \(A_r\).  Hence
\begin{equation}
 \QZ(\mathcal E_r)
 =2^{-(R-r)}\frac{A_r}{\cD_0}.
 \label{eq:tpc44-face-event-mass}
\end{equation}
The event \(\mathcal E_r\cap\{\varepsilon_{p_r}=-1\}\) is
\(\mathcal E_{r-1}\).  Dividing the corresponding masses gives
\begin{equation}
 \QZ(\varepsilon_{p_r}=-1\mid\mathcal E_r)
 =\frac{A_{r-1}}{2A_r}=\frac{1+\delta_r}{2}.
 \label{eq:tpc44-face-conditional-ratio}
\end{equation}
The product formula follows.
\end{proof}

If \(Z_0(-\one)=0\), every upper endpoint is automatic.  Otherwise, for
a target factor \(K\ge1\), exact endpoint selection is equivalent to
\begin{equation}
 \sum_{r=1}^{R}\log(1+\delta_r)\le\log K.
 \label{eq:tpc44-log-drift-criterion}
\end{equation}
This is why pointwise \(o(1)\) bounds for \(\delta_r\), or even a small
quadratic variation, need not suffice: a long sequence of tiny positive
drifts may still accumulate.
